\section{Introduction}\label{sec:intro}

Zirconium and its alloys --- Zircaloy-2 and Zircaloy-4 in light-water reactor
fuel cladding, Zr--2.5\,Nb in CANDU pressure tubes, and the low-tin variants M5
and ZIRLO in pressurized-water reactor service --- operate under prolonged
neutron exposure at displacement rates of $10^{-8}$ to
$10^{-7}\,\mathrm{dpa\,s^{-1}}$ and accumulate total damage of 5--30\,dpa over a
fuel-cycle lifetime~\citep{Adamson2019,LemaignanMotta1994}. The microstructural
changes that result --- nucleation and growth of \aloop{} prismatic loops, the
comparatively delayed appearance of \cloop{}-component basal vacancy loops,
dislocation channeling under load, anisotropic irradiation growth and creep, and
stress-driven hydride reorientation --- are the principal source of property
degradation in Zr-based components, manifested as hardening, loss of uniform
elongation, axial growth of fuel assemblies, and creep-down of pressure
tubes~\citep{Northwood1977,Griffiths1988,OnimusBechade2012,Cinbiz2016}.
Predicting those changes across the joint envelope of dose, temperature, texture
and applied stress is the central quantitative challenge of fuel-element and
pressure-tube qualification.

Cluster dynamics (CD), in the form it took in the late 1970s and early
1980s~\citep{Ghoniem1977,Ghoniem1979a,Ghoniem1980,Hayns1976,Kiritani1973},
remains the workhorse mesoscopic theory for that prediction. It treats each
defect species and size as a population whose concentration evolves under
coupled master rate equations balancing in-cascade production,
diffusion-controlled reactions, thermal emission and absorption at fixed
sinks~\citep{BrailsfordBullough1972,Wiedersich1972,Mansur1978}. Its advantages
are direct access to time scales far beyond molecular
dynamics~\citep{Bacon2009,Marian2017}, explicit resolution of the cluster-size
distribution and therefore of every quantity that enters a macroscopic property
model~\citep{GolubovSingh2001,Holt1988}, and modest cost relative to kinetic
Monte Carlo, which together make CD the meeting point between atomistic
parameterization~\citep{Christensen2015,PatraTrinkle2019} and engineering-scale
prediction~\citep{KohnertWirth2017,OnimusBechade2012}.

\subsection{The mean-field assumption, and where it fails}
\label{sec:intro-meanfield}

Classical CD is a \emph{mean-field, spatially averaged} theory. Every
concentration is a single number for the whole material, every sink is
represented by a sink strength $k^{2}$ that has already been averaged over an
assumed random and statistically uniform distribution of absorbers, and the
transport that connects a mobile defect to a sink appears only through that
average. Three distinct assumptions are folded into this, and each fails under
conditions that are not exotic.

\paragraph{Spatial uniformity.} A sink strength is meaningful only if the
diffusion field around each absorber has relaxed to a common far-field value
before it meets the next absorber. Near a free surface, a grain boundary, or a
denuded band, it has not: the mobile concentration varies by orders of magnitude
over a distance comparable with the boundary-layer width, and a single average
concentration is not merely imprecise, it does not describe any point in the
material. The consequences are structural rather than quantitative. In the
calculations reported below, the outermost shell of a
\SI{500}{\nano\meter} grain carries essentially the entire prismatic loop
population, so a domain average of the loop density is a statement about the
boundary layer and not about the material.

\paragraph{Diluteness.} The Avrami-type coalescence kernels that CD uses to
represent loop--loop and loop--network interaction are derived by assuming that
absorbers are placed independently and at random. Once the volume fraction swept
by growing loops becomes appreciable, the placement is no longer independent ---
loops that have already coalesced are correlated by construction --- and the
mean-field estimate of the overlap probability ceases to be an estimate of
anything. This is the point at which the population must be represented as
discrete objects rather than as a density, and locating it is a prediction the
model can make about itself.

\paragraph{Isotropy inherited from cubic metals.} Self-interstitial migration in
$\alpha$-Zr is strongly anisotropic, prism-plane mobility exceeding basal
mobility by roughly an order of magnitude, and vacancy migration is anisotropic
more weakly~\citep{Christensen2015,PatraTrinkle2019}. Most CD treatments of
zirconium nevertheless retain isotropic diffusion constants taken from the
body-centered cubic literature and recover the missing directionality through
independently fitted bias factors. That conflation is not a
book-keeping convenience: it allows a model to believe simultaneously that
diffusion is isotropic and that its sinks are biased by the anisotropy of that
same diffusion, and it severs the relation between atomic-scale parameters and
predicted dimensional change.

\subsection{Spatially resolved cluster dynamics}
\label{sec:intro-srcd}

Spatially resolved cluster dynamics (SRCD) replaces the averaged concentrations
by fields $c(\xv,t)$ obeying reaction--diffusion equations on a meshed domain,
so that the transport a mean-field sink strength summarizes is instead solved
for~\citep{spatialCD2024,christien2009cluster,barashev2015theoretical}. Three
things follow that a mean-field model cannot supply.

First, the boundary layer is resolved rather than parameterized. A grain
boundary in mean-field CD is a number, $k^2_{gb}\sim 6\sqrt{k^2}/d$; in SRCD it
is a Dirichlet condition, and the depleted zone in front of it --- whose width
is set by the competition between the local diffusivity and the local sink
density, and which therefore differs for each mobile species and evolves through
the transient --- is an output. The denuded zones seen in irradiated
microstructures are a direct consequence, and they are among the few
observations that speak to the transport parameters rather than to the reaction
parameters.

Second, quantities that are degenerate in a mean-field model become
distinguishable. \Cref{sec:coexistence} shows that diffusional anisotropy
difference opens a window of the mobile flux ratio over which prismatic
interstitial and basal vacancy loops grow together, but that the window is a
condition on \emph{one number}; a mean-field model carries that number once,
while a spatially resolved model carries it as a field, and the observed
coexistence of both loop characters can then be a statement about different
points of the same grain rather than about a single balance holding everywhere.

Third, the fields are the natural input to a discrete model. Dislocation
dynamics, micromechanical models of channeling, and fracture models all require
defect positions and sizes, not densities. A spatially resolved continuum field
carries the information a discrete population must reproduce --- local number,
local size, local orientation --- and the conversion between the two becomes a
well-posed problem with conservation laws attached, rather than a rescaling of a
single average~\citep{li2025coupledCDDD,ji2022concurrent,li2020coupled}.

\subsection{Self-consistency}
\label{sec:intro-selfconsistency}

\textbf{The model is self-consistent in two distinct senses, and both are
restrictive.} They are independent requirements --- a model can satisfy either
without the other --- and the term is used in this paper only for the
conjunction.

\paragraph{(a) Within the continuum model: every bias factor follows from
anisotropic diffusion theory.} Every capture efficiency that appears in a loop
growth law, and every sink strength that the same loops impose on the mobile
species, is generated from the \emph{same} anisotropic diffusion tensors that
transport those species, through the tensorial form of the
diffusional-anisotropy-difference (DAD)
correction~\citep{woo1988theory,woo1981sink,woo1982intrinsic,po2026tensorialdad}.
This holds for \emph{all} variants --- the three prismatic habits of each
character and the three basal states alike --- so one anisotropy parameter per
mobile species fixes the bias of every one of the nine families, and there is no
independent set of phenomenological bias factors anywhere in the model. What the
constraint rules out is worth stating as a possibility rather than an
abstraction: a model that diffuses the mobile species with one anisotropy while
biasing its sinks with an independently fitted pair can hold simultaneously the
two beliefs that diffusion is isotropic and that the sinks are biased by the
anisotropy of that diffusion. Generating both from one $p^{m}$ removes that
possibility by construction. The immediate consequence is that a change in the
transport anisotropy now moves the growth laws and the sink strengths together,
in the directions the transport dictates, and cannot be absorbed by refitting a
bias.

\paragraph{(b) Between the continuum and the discrete descriptions: the transfer
preserves what defines the microstructure.} The same population must be
representable as continuum fields or as a set of discrete loops, and the map
between the two must not create or destroy the quantities that make it the same
population. \Cref{sec:conversion} states the requirements and their ranking:
stored defects exactly, loop number to the integer draw, the dispersion of the
size distribution to the truncation of the sampling, and the sink strength
reported rather than forced. Sense~(b) is what makes a hybrid calculation
meaningful --- continuum where the population is dilute, discrete where it is
not, with the choice made on grounds of dose and computational cost rather than
of physics --- and it is what turns the conversion into a problem with
conservation laws attached instead of a rescaling of an average.

The two senses meet at the loop's capture cross-section. Sense~(a) makes that
cross-section a function of the transport tensors alone; sense~(b) requires the
discrete population to present the same cross-section as the continuum field it
replaces. A failure of either shows up in the transfer ledger of
\cref{sec:ledger} as a discontinuity in sink strength, which is why that ledger
is reported and not merely checked.

The immobile population is correspondingly expanded, from the lumped
stress-aligned pair of the earlier mean-field model to nine crystallographically
resolved families: the three prismatic variants $a_1,a_2,a_3$ of both
interstitial and vacancy character, and a three-state basal vacancy chain
running from the stacking-fault pyramid through the faulted loop to the perfect
loop. The expansion follows the experimental loop
character~\citep{jostsons1977nature,griffiths1988review,carpenter1981study,%
tournadre2012experimental,harte2017effect}: prismatic \aloop{} loops of both
signs, and basal \cloop{}-component loops that are predominantly \emph{faulted}
and that form by collapse of a pyramidal precursor rather than by accretion.
Each family carries the first three moments of its size distribution, so that
the width of the distribution --- not only its mean --- is transported, and the
conversion to a discrete population can draw loop sizes from a distribution
instead of assigning every loop at a point the same size.

\subsection{Contributions and outline}
\label{sec:intro-outline}

This paper reports the formulation, its numerical solution, and its first
application to a spatially resolved single crystal. Specifically:

\begin{enumerate}[label=(\roman*),itemsep=2pt]
\item an SRCD formulation for $\alpha$-Zr that is self-consistent in sense~(a):
      the capture efficiencies and sink strengths of all nine families are
      generated from the diffusion tensors, with no independent bias
      (\cref{sec:classes,sec:formulation});
\item a two-time-scale operator split that makes of order $10^{6}$ stiff
      ordinary differential equations per substep tractable, with the cost
      argument that justifies it and the splitting error that pays for it
      (\cref{sec:split});
\item a continuum-to-discrete conversion that is self-consistent in sense~(b),
      with its conservation guarantees stated and ranked and its residual
      discrepancies measured, and a grain-boundary treatment that makes the
      boundary a sink for immobile loops as well as for mobile defects
      (\cref{sec:conversion});
\item a calibration against 78 experimental measurements of loop density and
      diameter, presented as a maximum \emph{a posteriori} estimate under stated
      priors, with an explicit account of which parameters the data determine
      and which the bounds determine (\cref{sec:calibration}); and
\item a \SI{500}{\nano\meter} hexagonal single-crystal calculation to 10\,dpa
      reporting mobile and immobile fields, the closed atom balance including
      the boundary channel, the predicted denuded zone, and the discrete loop
      population the fields imply (\cref{sec:results}).
\end{enumerate}

\Cref{sec:classes} fixes the defect classes and the reactions among them.
\Cref{sec:formulation} gives the formulation, beginning from the general master
equation and reducing it to the equations actually integrated.
\Cref{sec:split} is the solution method. \Cref{sec:conversion} is the
continuum-to-discrete map and the grain-boundary models. \Cref{sec:calibration}
is the calibration. \Cref{sec:results} is the single-crystal study, and
\cref{sec:conclusions} draws the conclusions and states what remains open.
Four appendices collect the capture-efficiency derivation, the moment closure,
the peripheral-emission model that replaces fitted annealing lifetimes, and the
state-vector layout.
